\label{sec:background}

\subsection{Defining an Open Seat}

We use a strict and a broad definition of open seat. Under the \textit{strict} definition,
a district is open if no incumbent's home census tract is assigned to it:

\begin{definition}[Strict Open Seat]
District $d$ is a \emph{strict open seat} under plan $P$ if and only if:
\[
  \nexists\, i \in \mathcal{I} : d_P(t(i)) = d,
\]
where $\mathcal{I}$ is the set of incumbents seeking reelection, $t(i)$ is incumbent
$i$'s home census tract, and $d_P(t)$ is the district assignment of tract $t$ under $P$.
\end{definition}

Under the \textit{broad} definition, we additionally count districts that are
technically incumbent-occupied but where the incumbent faces a genuinely competitive
election due to redistricting-induced partisan composition change:

\begin{definition}[Broad Open Seat]
District $d$ is a \emph{broad open seat} under plan $P$ if either (a) it is a strict
open seat, or (b) it contains exactly one incumbent's home tract and its partisan lean
$v_d^D$ falls in the range $[0.45, 0.55]$ (competitive by the 5-percentage-point
threshold used in I.2).
\end{definition}

The broad definition captures the reality that a nominal incumbent district can
function like an open seat if redistricting has moved the partisan composition to
truly competitive territory---the incumbent will face challengers as well-funded and
experienced as those who contest open seats.

\subsection{Retirement Incentive Model}

We model the probability that an incumbent retires following redistricting as a
function of two factors: (a) whether she is paired with another incumbent (in which
case she faces a primary against a sitting colleague), and (b) how much her district's
partisan lean changed relative to her prior district. Formally, the retirement
incentive score for incumbent $i$ under plan $P$ is:

\[
  \text{RI}(i, P) = \alpha \cdot \mathbf{1}[\text{paired}(i,P)]
  + \beta \cdot |\Delta v^D_{d(i)}|,
\]

where $\Delta v^D_{d(i)} = v^D_{d_P(t(i))} - v^D_{d_\text{prior}(i)}$ is the change
in partisan lean of the district containing incumbent $i$'s home tract between the
prior map and the new plan $P$, and $\mathbf{1}[\text{paired}]$ is 1 if $i$ is in a
district with at least one other incumbent.

We set $\alpha = 0.40$ and $\beta = 0.50$ based on calibration against historical
post-redistricting retirement rates in cycles where pairing and lean-change data are
available~\citep{ansolabehere2000incumbency}. These coefficients are approximate; the
model is intended to illustrate the directional effect of algorithmic vs.\ enacted
redistricting on retirement pressure, not to produce precise predictions.

\subsection{Prior Literature on Open Seats and Redistricting}

The relationship between redistricting and open seats has been studied most thoroughly
in the context of the 1990 and 2000 redistricting cycles. \citet{ansolabehere2000incumbency}
find that redistricting increased open-seat rates in the 1990 cycle by approximately
30\% in states with major boundary changes. \citet{altman2011public} argue that public
mapping exercises consistently produce more competitive and more open maps than
legislative redistricting, consistent with the algorithmic comparison developed here.
\citet{king1994electoral} document that the electoral cycle effect on seat swings is
mediated by the open-seat rate: years with more open seats produce larger partisan swings.
